\documentclass[11pt,a4paper]{article}
\usepackage[T1]{fontenc}
\usepackage[utf8]{inputenc}
\usepackage{amsmath,amssymb,xcolor}
\usepackage[margin=25mm]{geometry}
\usepackage[colorlinks=true,urlcolor=blue]{hyperref}
\definecolor{nexorablue}{HTML}{2459B5}
\setlength{\parindent}{0pt}
\setlength{\parskip}{7pt}
\setlength{\emergencystretch}{2em}
\renewcommand{\familydefault}{\sfdefault}
\begin{document}

{\small\bfseries\color{nexorablue} NEXORA CPGE}\par
{\small MP / NEX-STR-05 / Corrigé}\par
{\LARGE\bfseries Un générateur explicite dans un anneau de polynômes}\par
\vspace{0.25em}\hrule\vspace{0.6em}

\section*{Divisibilité}
Chaque générateur est divisible par $X-1$, donc $I\subset(X-1)\mathbb R[X]$.

\section*{Identité et appartenance}
On calcule $(X^3-1)-X(X^2-1)=X-1$. Le membre de gauche appartient à $I$, qui contient donc tous les multiples de $X-1$. Finalement $P\in I$ équivaut à $P(1)=0$.


\end{document}
